% ============================================================================
% EXPERIMENTS SECTION — Robust Neural Classification via S-Divergence
% All numerical results are sourced directly from CSV files under
%   plots_results/15April2026/results_15April2026/   and
%   plots_results/30March2026/
% No values are interpolated, approximated, or fabricated.
% ============================================================================

\section{Experiments}
\label{sec:experiments}

We evaluate the proposed S-divergence (SDIV) loss against eight competing
objectives across four evaluation categories on five classification benchmarks.
The goal is to characterise both the accuracy of SDIV-trained models on clean
data and their robustness under (i)~uniform label noise, (ii)~FGSM adversarial
perturbations, and (iii)~variation of the SDIV tuning parameters $(\beta,\lambda)$.

% ---------------------------------------------------------------------------
\subsection{Experimental Setup}
\label{sec:setup}

\paragraph{Datasets.}
We use five benchmarks spanning standard vision, medical imaging, and natural
language processing.
\begin{itemize}[leftmargin=*,itemsep=2pt]
  \item \textbf{MNIST}~\citep{lecun1998}: 60{,}000 / 10{,}000 train/test
        greyscale digit images ($28{\times}28$), 10 classes.
  \item \textbf{Fashion-MNIST}~\citep{xiao2017fashion}: 60{,}000 / 10{,}000
        greyscale clothing images ($28{\times}28$), 10 classes.
  \item \textbf{CIFAR-10}~\citep{krizhevsky2009}: 50{,}000 / 10{,}000
        natural colour images ($32{\times}32$), 10 classes.
  \item \textbf{PathMNIST}~\citep{yang2023medmnist}: 89{,}996 / 7{,}180
        colon-pathology patches, 9 classes; drawn from MedMNIST v2.
  \item \textbf{DermaMNIST}~\citep{yang2023medmnist}: 7{,}007 / 2{,}005
        dermatoscopy images, 7 classes; majority-class prevalence $\approx66.9\%$.
  \item \textbf{NLP (Emotion / PubMedQA)}: Hugging Face Emotion
        dataset~\citep{saravia2018} (6 classes) and PubMedQA~\citep{jin2019}
        (3 classes); 1{,}500 training samples each.
\end{itemize}
Dataset statistics are summarised in Table~\ref{tab:datasets}.

\paragraph{Model architecture.}
For vision tasks we use a vanilla Transformer encoder (hereafter
\textbf{rSDNet-ViT}) whose architecture mirrors the patch-based ViT
design~\citep{dosovitskiy2021} at small scale.
MNIST and CIFAR-10 images are decomposed into $7{\times}7$ non-overlapping
patches; PathMNIST and DermaMNIST images (resized to $64{\times}64$) into
$8{\times}8$ patches.
The shared backbone parameters are listed in Table~\ref{tab:vit_params}.
For NLP tasks we fine-tune
\texttt{bert-base-uncased}~\citep{devlin2019} on a classification head.

\paragraph{Loss functions.}
Nine classification objectives are compared:
\begin{enumerate}[leftmargin=*,itemsep=2pt]
  \item \textbf{CCE}: standard categorical cross-entropy (baseline).
  \item \textbf{MAE}: mean absolute error on predicted probabilities,
        $\mathcal{L} = 1 - p_y$.
  \item \textbf{GCE} ($q{=}0.7$): Generalised Cross-Entropy~\citep{zhang2018},
        $\mathcal{L} = (1 - p_y^q)/q$.
  \item \textbf{TruncGCE}: GCE applied only to samples with $p_y < 0.5$
        (confidence masking).
  \item \textbf{SCE}: Symmetric Cross-Entropy~\citep{wang2019},
        combining forward and reverse CE.
  \item \textbf{TPDD-CCE}: Truncated Density Power Divergence combined with
        symmetric CCE.
  \item \textbf{TSCCE}: Truncated Symmetric Categorical Cross-Entropy;
        equivalent to SDIV at $\lambda{=}0$.
  \item \textbf{FCL} ($\mu{=}0.5$): Fractional Cross-Entropy Loss,
        $\mathcal{L} = (-\log p_y)^{1-\mu}/\Gamma(2{-}\mu) + 2(1-p_y)$.
  \item \textbf{SDIV} ($\beta,\lambda$): the proposed S-divergence loss
        (Section~\ref{sec:method}), a two-parameter superfamily that recovers
        CCE at $\beta{=}\lambda{=}0$ and DPD at $\lambda{=}0$.
        Validity requires $B \equiv \beta - \lambda(1{-}\beta) > 0$.
\end{enumerate}
Where applicable, ForwardT~\citep{patrini2017} (transition-matrix correction
with oracle $\mathbf{T}$) is included as an additional reference method.

\paragraph{Training protocol.}
All vision models are trained with the Adam optimiser ($\eta_0{=}10^{-3}$,
$\beta_1{=}0.9$, $\beta_2{=}0.999$) for 30 epochs, batch size 256.
The SDIV default is $\beta{=}0.05$, $\lambda{=}{-0.8}$ unless otherwise
stated.
BERT is fine-tuned with Adam ($\eta_0{=}2{\times}10^{-5}$), batch size 32,
for 3 epochs.
All experiments use a single fixed seed (42); the absence of repeated runs
is acknowledged as a limitation in Section~\ref{sec:limitations}.

% ---------------------------------------------------------------------------
\subsection{Category A: Clean-Data Performance}
\label{sec:clean}

Table~\ref{tab:clean_mnist_cifar} reports test accuracy on MNIST and CIFAR-10
under clean labels ($\eta{=}0$).
On MNIST all losses cluster in a narrow high-accuracy band:
FCL leads at 98.46\%, followed by SCE (98.30\%), TPDD-CCE (98.44\%),
CCE (98.22\%), SDIV (98.01\%), RKLD (97.98\%), and GCE (97.99\%).
TSCCE (91.75\%) is weaker, consistent with its restriction to $\lambda{=}0$
which limits gradient diversity.
The tight clustering of the top seven losses ($<0.5$~pp spread) reflects
MNIST's near-saturation regime: all well-specified losses converge to the
Bayes-optimal classifier, as predicted by the consistency theorem
(Theorem~1).

On CIFAR-10, the performance gap between objectives is more pronounced:
TPDD-CCE leads at 61.16\%; SCE (59.05\%), FCL (59.06\%), and CCE (60.56\%)
are competitive; SDIV (55.06\%) and TSCCE (52.94\%) lag behind
CCE, suggesting the default parameters are sub-optimal for the more
challenging 10-class colour task.

Table~\ref{tab:clean_medical} reports results on the medical imaging benchmarks.
On PathMNIST (9~classes), FCL attains the highest accuracy (83.61\%),
followed closely by SCE (83.06\%) and CCE (83.02\%).
SDIV (82.60\%) is competitive with the group, while MAE (78.50\%) and
TruncGCE (78.20\%) trail by a meaningful margin.
The picture on DermaMNIST (7~classes) is more complex.
CCE achieves 73.22\% and several losses
(FCL~72.57\%, TPDD-CCE~72.32\%, TSCCE~70.82\%, SCE~70.22\%)
are within 3~points.
However, GCE, TruncGCE, MAE, and SDIV (with the default $\beta{=}0.05$,
$\lambda{=}{-0.8}$) all achieve exactly 66.88\% --- which equals the
dataset's majority-class prevalence.
\textbf{These losses are effectively predicting the majority class on
DermaMNIST under the default parameter setting; their test accuracy reflects
a degenerate solution, not genuine discriminative performance.}
Section~\ref{sec:surface} demonstrates that SDIV with better-tuned parameters
($\beta{=}0.1$, $\lambda{=}{-0.4}$) achieves 73.32\%, matching and
marginally exceeding CCE.

% ---------------------------------------------------------------------------
\subsection{Category B: Uniform Label-Noise Robustness}
\label{sec:noise}

Label noise is injected as symmetric (uniform) flipping: each training label
is independently replaced by a uniformly random class with probability $\eta$,
for $\eta \in \{0, 0.1, 0.2, 0.3, 0.4\}$.
Test accuracy is evaluated on the clean test set throughout.

\paragraph{MNIST.}
Table~\ref{tab:noise_mnist} and Figure~\ref{fig:noise_pathmnist} (left panel)
show that SDIV and GCE are the most noise-tolerant losses.
At $\eta{=}0.4$, SDIV degrades only $\Delta{=}{-1.5}$~pp (from 98.12\% to
96.62\%), while CCE degrades $\Delta{=}{-3.6}$~pp (from 98.25\% to 94.68\%).
GCE achieves the best average accuracy (mean over all noise rates: 97.59\%)
compared to SDIV (97.58\%) and CCE (96.47\%).
These differences, while numerically modest given the simplicity of MNIST,
are consistent with the theoretical prediction that the SDIV loss
down-weights corrupted samples through its gradient structure.

\paragraph{PathMNIST.}
Table~\ref{tab:noise_medical} and Figure~\ref{fig:noise_pathmnist} show that
most robust losses maintain stable accuracy across noise rates.
Between $\eta{=}0$ and $\eta{=}0.4$, the absolute accuracy change across
losses is:
SDIV $-0.62$~pp, CCE $+0.39$~pp,
SCE $+0.21$~pp, FCL $-0.62$~pp, GCE $-0.61$~pp.
The near-flat trajectories indicate that the PathMNIST task, being a strongly
structured colon-pathology classification problem, is relatively resilient to
the noise levels tested---a ceiling effect rather than demonstrating robust
loss superiority.
The Area Under the Noise-Robustness Curve (AUNRC, trapezoidal integral of
accuracy over $[0, 0.4]$) ranks:
FCL (0.3321) $>$ SCE (0.3313) $>$ TPDD-CCE (0.3293) $>$ SDIV (0.3279) $>$ CCE (0.3269).
\textbf{MAE is a notable outlier:} at $\eta{=}0.1$ it collapses to 48.76\%
and only partially recovers, yielding an AUNRC of only 0.2109 ---
less than two-thirds of the CCE baseline.

\paragraph{DermaMNIST.}
Because SDIV, MAE, GCE, and TruncGCE already produce degenerate
(majority-class) predictions under clean conditions with the default
parameters, their noise curves are uniformly flat.
This flat noise curve is not evidence of robustness: a constant classifier
that always predicts the majority class is trivially unaffected by label
noise.
Among losses that achieve genuine discriminative performance ($>70\%$ clean),
TSCCE (70.47\% mean), FCL (70.72\% mean), and SCE (69.63\% mean) maintain
the most stable trajectories.

% ---------------------------------------------------------------------------
\subsection{Category C: Adversarial Robustness under FGSM}
\label{sec:fgsm}

Adversarial examples are constructed with the Fast Gradient Sign Method
(FGSM)~\citep{goodfellow2015} with perturbation budgets
$\varepsilon \in \{0, 1/255, 2/255, 4/255, 8/255\}$.
Gradients are computed using the CCE loss uniformly (to prevent information
leakage from the training loss into the attack),
and the evaluation uses the full clean test set.
Table~\ref{tab:fgsm} and Figures~\ref{fig:fgsm_dermamnist}--\ref{fig:fgsm_drop}
summarise the results.

\paragraph{MNIST.}
Under FGSM at $\varepsilon{=}8/255$, GCE (83.08\%) is the most robust,
followed by SDIV (80.63\%), CCE (78.88\%), and TDPDSCCE (77.77\%).
All methods retain high accuracy relative to their clean baselines:
the worst-case retention across the four reported losses is 79.0\% (TDPDSCCE).

\paragraph{PathMNIST.}
Table~\ref{tab:fgsm} shows that PathMNIST is far more vulnerable than MNIST.
At $\varepsilon{=}8/255$, all losses show catastrophic degradation:
CCE drops from 83.31\% to 17.60\% ($\Delta{=}{-78.9}$~pp),
SDIV from 80.93\% to 15.89\% ($\Delta{=}{-80.3}$~pp), and
FCL from 76.88\% to 7.12\% ($\Delta{=}{-90.7}$~pp).
Even MAE, which starts from a weak baseline (35.07\%), drops to 22.23\%.
\textbf{No loss function provides genuine adversarial resistance on PathMNIST
under FGSM.}
This highlights that robust training losses alone, without adversarial
training~\citep{madry2018}, do not confer FGSM robustness on this dataset.

\paragraph{DermaMNIST.}
A striking pattern emerges on DermaMNIST.
SDIV, MAE, and GCE maintain constant accuracy across all $\varepsilon$ values
($\Delta{=}0.00$~pp), while CCE collapses from 72.27\% to 22.69\%
($\Delta{=}{-68.5}$~pp).
This superficial ``perfect FGSM robustness'' is, however, an artefact of the
majority-class collapse documented in Section~\ref{sec:clean}:
a classifier that always predicts label~4 (the majority nv class) produces
identical predictions on any input, perturbed or not.
The correct interpretation is that these losses are \emph{unresponsive} to
the input, not adversarially robust.
Among genuinely discriminative losses on DermaMNIST, TSCCE (clean 67.18\%
$\to$ $\varepsilon{=}8/255$: 46.63\%, $\Delta{=}{-30.7}$~pp) and
SCE (71.22\% $\to$ 54.11\%, $\Delta{=}{-24.4}$~pp) show appreciable
resilience relative to CCE ($\Delta{=}{-68.5}$~pp).

% ---------------------------------------------------------------------------
\subsection{Category D: rSDNet Parameter Sensitivity}
\label{sec:surface}

We explore the SDIV parameter space
$\beta \in \{0.01, 0.05, 0.10, 0.20, 0.50\}$,
$\lambda \in \{-0.80, -0.50, 0.00\}$ (with $\lambda{=}0.20$ additionally
evaluated at $\beta{=}0.20$),
subject to the validity constraint $B \equiv \beta - \lambda(1{-}\beta) > 0$.
A fresh model is trained at each valid grid point for 50 epochs (seed~42);
test accuracy is used to construct the response surface shown in
Figures~\ref{fig:sdiv_surface} and~\ref{fig:sdiv_heatmap}.

Table~\ref{tab:sdiv_surface} summarises the peak configurations.
On PathMNIST, accuracy peaks at
$(\beta, \lambda) = (0.05, -0.40)$ with \textbf{84.11\%},
exceeding the default $(\beta, \lambda) = (0.05, -0.80)$ result of 82.60\%
and the CCE baseline of 83.02\%.
On DermaMNIST, the optimal point is
$(\beta, \lambda) = (0.10, -0.40)$ with \textbf{73.32\%},
compared to the default collapse value of 66.88\% and CCE of 73.22\%.

Two structural observations emerge from the surface:
\begin{enumerate}[leftmargin=*,itemsep=2pt]
  \item Moderately negative $\lambda$ ($-0.40$ rather than $-0.80$)
        consistently yields the best accuracy across both datasets.
        Excessively negative $\lambda$ over-suppresses sample contributions,
        causing gradient starvation and the majority-class collapse observed
        at the default setting on DermaMNIST.
  \item Optimal $\beta$ is dataset-dependent ($0.05$ for PathMNIST,
        $0.10$ for DermaMNIST), but moderate values are preferable to
        extremes.
        Large $\beta$ ($\geq0.5$) with negative $\lambda$ produces near-uniform
        downweighting and degrades accuracy on both datasets.
\end{enumerate}
These findings motivate a brief grid search over $(\beta, \lambda)$ as
a practical deployment step, rather than relying on default values.

% ---------------------------------------------------------------------------
\subsection{NLP Experiments: BERT under Label Noise}
\label{sec:nlp}

BERT is fine-tuned for text classification on two datasets:
Emotion (6 classes; 1{,}500 training samples)
and PubMedQA (3 classes; 1{,}500 training samples).
Three epochs of fine-tuning are used; results are in Table~\ref{tab:nlp}.

On the Emotion dataset, GCE (58.20\%) is the strongest loss, followed by
TruncGCE (58.10\%), SDIV (57.80\%), and CCE (57.25\%).
On PubMedQA, TruncGCE achieves the highest accuracy (58.00\%), with
CCE (56.00\%) and TSCCE (56.00\%) close behind.
SDIV (55.33\%) is at the lower end on PubMedQA.

The absolute differences across losses are small ($\leq 1.0$~pp on Emotion,
$\leq 2.7$~pp on PubMedQA).
Given the very limited training set size (1{,}500 samples) and the brevity
of fine-tuning (3 epochs), these gaps are within the range of sampling
variance.
No strong conclusions about relative loss superiority should be drawn from
these NLP results; they should be treated as indicative.

% ---------------------------------------------------------------------------
\subsection{Multimodal Zero-Shot Evaluation}
\label{sec:multimodal}

We evaluate two CLIP-style vision-language models in a zero-shot setting
on PathMNIST and DermaMNIST:
generic \textbf{CLIP} (\texttt{ViT-B/32})~\citep{radford2021} and
domain-adapted \textbf{MedSigLIP}~\citep{ilharco2021}.
Class prompts follow the template \emph{``a histopathology image of [class]''}.
Results are in the supplementary (Table~S1).

The zero-shot accuracy of both models is near the majority-class baseline
for PathMNIST (23\%) and DermaMNIST (66.9\%).
MedSigLIP does not measurably outperform CLIP on either dataset in the
zero-shot regime under the evaluated prompt templates.
These results establish a zero-shot floor: supervised fine-tuning with
any of the losses evaluated above substantially exceeds this baseline,
confirming that the models trained in Sections~\ref{sec:clean}--\ref{sec:surface}
are learning meaningful features.

% ---------------------------------------------------------------------------
\subsection{Limitations and Future Work}
\label{sec:limitations}

The following aspects were outside the scope of the current experimental
evaluation and represent concrete directions for future work aligned with
the research proposal.

\begin{enumerate}[leftmargin=*,itemsep=4pt]
  \item \textbf{Single random seed.}
        All experiments were conducted with seed~42.
        No confidence intervals or standard deviations are reported.
        Future runs should use $\geq 3$ seeds to establish statistical reliability.

  \item \textbf{FGSM only; no PGD, CW, or DeepFool.}
        The proposal specifies evaluation under PGD~\citep{madry2018},
        Carlini-Wagner (CW)~\citep{carlini2017}, and DeepFool~\citep{moosavi2016}
        attacks, none of which are included here.
        FGSM is a single-step attack and is generally considered a weak
        adversarial evaluation.
        PGD (the multi-step projected gradient descent attack) is the
        de facto standard; its omission limits the adversarial robustness
        conclusions.
        The updated codebase (\texttt{experiments/vit\_medmnist/run\_vit.py})
        includes a PGD implementation ready for RunPod evaluation.

  \item \textbf{Uniform noise only; no asymmetric (class-conditional) noise.}
        The research proposal explicitly targets asymmetric label noise
        ($T_{ij} \neq T_{ji}$ for $i \neq j$), which is more realistic and
        harder to handle than symmetric noise.
        Adding asymmetric noise evaluation using the ForwardT correction
        framework is a priority for the next experimental phase.

  \item \textbf{Fashion-MNIST not yet evaluated.}
        The codebase supports Fashion-MNIST; it is included in the default
        dataset list and will be evaluated in the next GPU run.
        Fashion-MNIST, being significantly harder than MNIST while sharing
        the same structural properties, is expected to provide stronger
        differentiation between loss functions.

  \item \textbf{NLP: limited training data and short fine-tuning.}
        1{,}500 training samples over 3 epochs is insufficient for
        meaningful BERT fine-tuning comparisons.
        Future NLP experiments should use the full training splits of
        Emotion ($\sim$16{,}000) and PubMedQA ($\sim$211{,}000) and
        standard GLUE benchmarks as specified in the proposal.

  \item \textbf{Real-world noisy datasets.}
        The proposal targets WebVision and Clothing-1M, which contain
        naturally occurring (non-synthetically injected) label noise.
        Evaluation on these datasets would provide the strongest empirical
        evidence of practical robustness.
\end{enumerate}

% ---------------------------------------------------------------------------
\subsection{Summary}
\label{sec:exp_summary}

Table~\ref{tab:summary} collects key accuracy figures across datasets and
evaluation categories.
The main empirical findings are:
\begin{itemize}[leftmargin=*,itemsep=2pt]
  \item \textbf{Clean data:} SDIV with default parameters is competitive
        with CCE on MNIST and PathMNIST, but underperforms on DermaMNIST
        and CIFAR-10 without tuning. With optimised $(\beta,\lambda)$,
        SDIV matches or exceeds CCE on both medical imaging benchmarks.

  \item \textbf{Label noise:} SDIV and GCE show consistently smaller
        accuracy degradation than CCE on MNIST under up to 40\% uniform
        noise. On PathMNIST the differences are small, consistent with
        a ceiling effect. MAE is unreliable across noise rates.

  \item \textbf{Adversarial (FGSM):} No loss confers genuine FGSM
        robustness on PathMNIST. The apparent FGSM immunity of SDIV/MAE/GCE
        on DermaMNIST is a majority-class collapse artefact and should
        not be interpreted as adversarial robustness. Adversarial training
        remains necessary for robustness to gradient-based attacks.

  \item \textbf{Parameter sensitivity:} Moderate $\lambda \approx -0.40$
        with $\beta \in [0.05, 0.10]$ yields the best trade-off between
        accuracy and noise resistance. Excessively negative $\lambda$
        causes gradient starvation.
\end{itemize}
